\chapter{研究目标、内容与关键问题}\label{chap:content}

\section{研究目标}

本课题面向复杂度尚未确定的 NP 判定问题，研究如何引入人工智能辅助构造多项式时间归约，并在选定问题上尝试证明其 NP 完全性（以下简称 NPC）。人工智能在该过程中不替代数学证明，而用于组织源问题检索、候选映射生成、实例反例搜索和候选修正。最终结论仍需写成可核验的复杂性证明。

设源问题和目标问题对应的语言分别为

\begin{equation}
    L_1,L_2\subseteq\Sigma^*,
    \label{eq:source-target-languages}
\end{equation}

其中 \(L_1\) 是已知 NP 困难问题，\(L_2\) 是待研究的问题。候选归约由函数 \(f:\Sigma^*\rightarrow\Sigma^*\) 给出。若要证明 \(L_2\) 的 NP 困难性，需要证明

\begin{equation}
    \forall x\in\Sigma^*,\qquad
    x\in L_1 \Longleftrightarrow f(x)\in L_2,
    \label{eq:target-specification}
\end{equation}

并给出 \(f\) 的多项式运行时间和输出长度上界。若还能证明 \(L_2\in\mathrm{NP}\)，则可得到 \(L_2\) 的 NP 完全性。实例检验只用于发现不满足式~\eqref{eq:target-specification}的输入，不能代替对所有输入的证明。

\section{研究内容}

现有研究尚未形成从可执行构造到未知问题复杂性证明的完整路径。自然语言归约数据集保存了证明文本，但评价仍依赖专家判断。领域专用语言和自动验证方法支持候选归约的自动检查，适用范围则受表示形式限制。NP 问题求解基准只检查单个实例的答案，通用代码智能体的测试结果也不直接对应归约正确性。复杂度未知的问题没有参考映射，还需处理源问题选择和新映射构造。

未知问题不存在可供对照的标准归约，若直接在该问题上开发方法，构造失败难以归因。错误既可能来自问题定义和判定程序，也可能来自映射构造或智能体的工具使用。为获得可控的研究对象，本课题先把已知 Karp 归约改写为程序构造任务。归约程序接收源问题实例并输出目标问题实例，判定程序在隐藏实例上检查两者的 YES/NO 等价性。可信参考归约用于确认任务存在正确实现，对抗性实例用于暴露阈值、索引和边界条件错误。该转换过程如图~\ref{fig:npc-reduction-coding}所示。数据集为归约构造提供可重复的实例级评价，也为后续分析失败原因保留统一材料。

\begin{figure}[!htbp]
    \centering
    \includegraphics[width=0.96\textwidth]{npc-reduction-coding}
    \bicaption{NPC 归约问题向程序构造任务的转换。}{Conversion of an NPC reduction problem into a coding task.}
    \label{fig:npc-reduction-coding}
\end{figure}

归约映射的程序化表达适合由代码智能体反复编辑和执行。本课题将基础代码智能体作为阶段性任务，开发任务工作空间、文件编辑、终端执行和多轮交互功能，并完成归约任务接入、运行控制和提交接口。基础代码智能体为数据集任务和后续专用工具提供统一运行环境，但通用程序反馈无法判断目标实例格式及两类问题的判定等价性。

本课题在基础代码智能体上开发 NPC 归约评测框架。该框架通过专用工具和技能模块提供问题定义、格式校验、判定反馈和反例信息。智能体根据这些信息修改归约映射，逐步排除输出格式错误、判定不一致和规模膨胀。有限实例检验只用于发现错误，不能代替归约证明。

已知归约任务用于检验候选生成与反例修正过程，复杂度未知的问题用于考察这套方法能否支持新的理论分析。课题将选择具有明确输入编码和 NP 成员性论证的目标问题，再从已知 NP 完全问题中寻找结构相近的源问题。候选程序与失败实例将用于提取映射规则和局部性质。最终结果需要给出双向蕴含及多项式资源界的数学证明。若完整证明尚未形成，则保留适用范围明确的困难性结论、局部引理或排除某类构造的反例。

\section{总体研究路线}

总体路线如图~\ref{fig:research-route}所示。已知归约数据集提供问题定义、程序入口和实例级评价。基础代码智能体负责归约任务接入、运行控制和归约程序执行。NPC 归约评测框架通过专用工具和技能模块补充格式校验、判定反馈与反例信息。未知复杂度问题案例用于检验这套方法能否辅助构造新的多项式时间归约。

\begin{figure}[!htbp]
    \centering
    \begin{tabular}{c}
        \fbox{\parbox{0.68\textwidth}{\centering
        NPC 归约程序构造数据集\\
        已知归约任务与实例级评价}}\\[0.35em]
        \(\Downarrow\)\\[0.35em]
        \fbox{\parbox{0.68\textwidth}{\centering
        基础代码智能体\\
        任务接入、运行控制与程序执行}}\\[0.35em]
        \(\Downarrow\)\\[0.35em]
        \fbox{\parbox{0.68\textwidth}{\centering
        NPC 归约评测框架\\
        专用工具、技能模块与判定反馈}}\\[0.35em]
        \(\Downarrow\)\\[0.35em]
        \fbox{\parbox{0.68\textwidth}{\centering
        未知复杂度问题案例\\
        源问题选择、候选构造与反例修正}}\\[0.35em]
        \(\Downarrow\)\\[0.35em]
        \fbox{\parbox{0.68\textwidth}{\centering
        NPC 证明或部分理论结论\\
        双向蕴含、资源界与适用范围}}
    \end{tabular}
    \bicaption{总体研究路线。}{Overall research route.}
    \label{fig:research-route}
\end{figure}
\FloatBarrier

数据集提供已知归约任务，基础代码智能体负责运行归约构造过程，评测框架提供归约相关反馈。归约程序通过有限实例检验后，还要改写为数学映射，并证明两个蕴含方向和多项式资源界。未知问题阶段不再以任务通过率作为最终评价，而以完整 NPC 证明或适用范围明确的部分理论结论为准。

\section{关键问题与技术难点}

源问题需要与目标问题具有可利用的结构对应关系。候选映射必须保持任意输入上的 YES/NO 答案，阈值、索引和边界情形都可能使归约失效。

判定程序与测试实例决定反馈是否可信。源问题和目标问题的判定结果必须准确，测试集需要覆盖 YES、NO 和典型边界实例，避免错误候选因测试不足而通过。

候选程序需要对应真正的多项式时间映射。有限测试无法排除指数搜索或先求解源实例再输出固定目标实例的做法，最终仍需证明双向蕴含、运行时间界和输出长度界。
